% SPDX-License-Identifier: Apache-2.0
% covers: Rom
\datasheet{Rom}{Vreteno's boot memory on the bus: the 1024 words the core fetches, readable and not writable, behind an AXI peripheral end.}

\dsfacts{\code{cpu/vreteno/src/rom.rs}}{\code{vreteno32::rom::Rom}}%
{\code{//docs:vreteno}}

\subsection*{Function}

\code{Rom} is the boot memory of the Vreteno machine as a device on
its bus: the same \code{IMEM\_WORDS}, 1024 words, that the core
fetches from inside itself, at address zero, so that a load can read
a constant that sits beside the code, and a debugger or a host on the
bus can read the program that is there. The core keeps its own copy
for the fetch, which reads it in the same cycle; this copy answers the
bus a cycle after a read is taken, as the data memory does. Both are
filled with the same image before the first cycle, so they cannot
disagree, and neither is written at run time.

It is the peripheral client of an \code{AxiPer}: it takes requests on
\code{req} and write beats on \code{wd}, and answers on \code{ans}
for a write and \code{rb} for a read. It answers single-beat bursts,
which is all the core makes.

\subsection*{Parameters}

\begin{center}\footnotesize
\begin{tabular}{@{}lp{0.7\textwidth}@{}}
\toprule
Parameter & Meaning \\
\midrule
\code{I} & The width of the AXI identifier on \code{req}, \code{ans}
and \code{rb}, in bits. The tables use 2, as the board does. \\
\bottomrule
\end{tabular}
\end{center}

\subsection*{Address map}

The word a burst names is bits 11 to 2 of its address. The memory
looks at no other address bits.

On the board, \code{BoardRouter} sends it the page at \code{0} with
mask \code{0xffff\_f000}, so addresses \code{0x0} to \code{0xffc} are
words 0 to 1023, which are the words the core fetches from
\code{0x0} to \code{0xffc}. \code{run.rs} puts it at the same place.

\dstables{Rom}

\subsection*{Behaviour}

\begin{itemize}
\item A read is taken when no write is waiting for its beat and the
answer register is free or is being emptied in this cycle. The word
is read into \code{word} at the edge, with its identifier in
\code{rid}, and \code{answer} is set.
\item The read is answered on \code{rb} the cycle after, with
\code{rid}, \code{OKAY} and \code{last} set.
\item A write is taken when no other write is waiting for its beat.
Its identifier goes into \code{pid} and \code{pend} is set; nothing
is written.
\item The held write's beat is taken when it has arrived and
\code{ans} has room, and dropped, so that the write data channel is
not left holding it; \code{ans} sends \code{SLVERR} with \code{pid},
which is what a peripheral that was reached and refused says.
\item While a write waits for its beat, no request of either kind is
taken.
\item \code{Rom::with} fills the words from a program before the
first cycle: the same words \code{Vreteno::with} puts in the core.
\end{itemize}

\subsection*{Verification}

\begin{itemize}
\item \code{the\_boot\_memory\_is\_readable\_and\_not\_writable}, in
\code{//cpu/vreteno:board\_test}, runs the board under a program
that loads a word of itself from address zero, stores over it, loads
it again and says whether it kept; the first load is checked to
have returned the program rather than the zero a hole answers.
\item \code{//cpu/vreteno:hello\_test},
\code{//cpu/vreteno:hello\_cc\_test} and
\code{//cpu/vreteno:traps\_test} run compiled programs whose
constants the linker put beside the code, in \code{.rodata}, which
their loads read through it.
\item Its netlist is part of the board's, which
\code{//cpu/vreteno:vreteno\_board\_synth} synthesises and the manual
\code{//cpu/vreteno/board/sim:board\_test} runs under Vivado's
simulator; it is not simulated on its own.
\end{itemize}

\subsection*{Used by}

\code{Board}, as \code{rom} behind the tracker \code{prom}, at the
router's first slot. \code{run.rs}, which every compiled program's
test runs on, as \code{Rom<IW>}.

\subsection*{Limits}

It serves bursts of one beat and does not look at a burst's length. It
checks no address, since the router sends it only its own range, so
an address beyond word 1023 wraps. One write is held at a time. It is
a second copy of the program, in block RAM beside the core's, and the
two are the same only because the netlist fills both from one image:
nothing at run time can change either.
